\chapter*{Glosario}
\phantomsection
\addcontentsline{toc}{chapter}{Glosario}
\markboth{Glosario}{Glosario}
\thispagestyle{empty}

\begin{description}
  \item[Bare-metal] Modo de ejecución en que el software corre directamente sobre
        el hardware sin sistema operativo intermedio. En este trabajo se refiere
        al firmware que se ejecuta sobre el CV32E40P dentro del SoC PULPissimo,
        accediendo a los periféricos y registros CSR sin ninguna capa de
        abstracción de sistema operativo.

  \item[APU] Unidad de procesamiento auxiliar (\textit{Auxiliary Processing
        Unit}). Interfaz del CV32E40P que enlaza el núcleo con coprocesadores
        externos como la FPU.

  \item[BRAM] Memoria embebida en bloque (\textit{Block RAM}). Recurso de
        memoria síncrona integrado en las FPGA, organizado en bloques de
        capacidad fija.

  \item[CV32E40P] Procesador RISC-V de 32 bits desarrollado por el OpenHW Group
        como evolución del núcleo RI5CY del equipo PULP. Implementa el ISA
        RV32IMFC con las extensiones Xpulp y constituye el núcleo central del
        SoC utilizado en este trabajo.

  \item[GDB] Depurador usado para cargar firmware, controlar la ejecución y leer
        registros del procesador a través de OpenOCD.

  \item[HPC] Cómputo de alto rendimiento (\textit{High-Performance Computing}).
        Categoría de plataformas con sistema operativo y contadores de hardware
        expuestos para telemetría energética estandarizada.

  \item[IoT] Internet de las cosas (\textit{Internet of Things}). Ecosistema de
        dispositivos embebidos conectados que operan con presupuestos de potencia
        limitados.

  \item[ISA] Arquitectura del conjunto de instrucciones
        (\textit{Instruction Set Architecture}). Define la interfaz visible entre
        el software y el procesador.

  \item[JTAG] Interfaz estándar IEEE~1149.1 usada para prueba y depuración de
        circuitos mediante señales serie como TCK, TMS, TDI y TDO.

  \item[LUT] Tabla de búsqueda (\textit{Look-Up Table}). Bloque combinacional
        básico de las FPGA, configurable para implementar cualquier función
        lógica de pocas entradas.

  \item[MPSSE] Motor síncrono multiprotocolo de FTDI
        (\textit{Multi-Protocol Synchronous Serial Engine}) usado para generar
        señales JTAG desde el FT232H.

  \item[OpenOCD] Herramienta de depuración que controla el acceso JTAG y expone
        una interfaz para que GDB pueda cargar programas y leer registros del
        procesador.

  \item[SIMD] Una instrucción, múltiples datos
        (\textit{Single Instruction Multiple Data}). Estilo de ejecución
        vectorial en el que una misma instrucción procesa varios operandos
        empaquetados en un mismo registro.

  \item[SoC] Sistema en un chip (\textit{System-on-Chip}). Integración de
        procesador, memoria, buses y periféricos dentro de una misma plataforma.

  \item[TCDM] Memoria de datos estrechamente acoplada
        (\textit{Tightly Coupled Data Memory}). En PULPissimo se refiere al acceso
        directo del procesador a la SRAM interna.
\end{description}

\cleardoublepage

%%% Local Variables:
%%% mode: latex
%%% TeX-master: "main"
%%% End:
